\odiseachapter{Menelao}{Menelao}\label{cap:menelao}

En el capítulo anterior, protestaba de la pobre lectura que nos ofrece Nolan de Helena. Eso no quiere decir ---he insistido ya en el tema en los capítulos precedentes--- que el afamado director estuviera obligado a ofrecer la misma versión de la divina que presenta Homero. De hecho, puede argumentarse que el único sentido que tiene revisitar las épicas clásicas es reinventárselas. Mi queja no es que Nolan reinterprete a Helena, sino que sale del paso con un cliché.

Permíteme, paciente lectora, generoso lector, que ilustre lo que quiero decir ---con toda desvergüenza--- usando un fragmento nuevo de \emph{Abandonando Ítaca}.

En la versión publicada del texto, Helena y Menelao aparecen solo de refilón. La escritura del capítulo anterior, sin embargo, me inspiró una extensión al poema, que ofrezco aquí a la generosa lectora, al paciente lector. Ulises empieza por recordar la (imaginaria) reunión anual, en Esparta, de los héroes que todavía quedan vivos, todos ellos ya muy ancianos.

\begin{versopropio}
Uno de los mejores momentos del año\\
es cuando viajas a Esparta.\\
Menelao organiza una reunión, a la que acuden\\
los que sobrevivieron a Troya y al \emph{nóstos}\footnote{El regreso a la patria después de la guerra.}.
\stanzabreak
Por supuesto, eres el invitado de honor,\\
te agasajan con comida y vino,\\
y hay una linda esclava de suaves manos\\
que cuida bien de tu cuerpo cansado.
\stanzabreak
Y todos disfrutáis, recordando\\
episodios repetidos tantas veces con los años\\
que es difícil separar la verdad de la fantasía,\\
la realidad de la imaginación.
\stanzabreak
A veces sospechas que la mitad de tus hazañas\\
no son más que reconstrucciones, tan lindas como tu esclava,\\
de algo más simple y sucio y brutal,\\
conjuradas para cubrir la sangre y la mugre y el estiércol.
\stanzabreak
¿Pero a quién le importa ya? Os reunís cada año,\\
y cada año hay menos y menos invitados en casa de Menelao,\\
y coméis, y bebéis, y cantáis, para celebrar\\
el fingimiento de que vuestras vidas tuvieron algún sentido.
\end{versopropio}

Tras lamentarse de que la mayoría de sus compañeros han muerto, Odiseo reflexiona sobre los dioses:

\begin{versopropio}
Cada vez sois menos. La mayoría ya no está.\\
Muchos de tus amigos no volvieron de Troya,\\
Aquiles, Áyax, Patroclo, Antíloco,\\
la lista es larga. Otros murieron en el viaje de vuelta.
\stanzabreak
Mucho se ha escrito sobre los dioses\\
oponiéndose al regreso de los argivos,\\
y poco sobre las ratas que, sospechas,\\
estaban detrás de las enfermedades que mataron a tantos.
\stanzabreak
Ah, pero a la gente le encanta cantar a los inmortales,\\
ruines y caprichosos como son,\\
y celebrar sus enredos en los asuntos humanos,\\
su lujuria y su crueldad dementes.
\stanzabreak
Puestos a elegir, tú te quedarías con las ratas,\\
que, al menos, iban a lo suyo,\\
y aunque se comían tu comida y se bebían tu agua,\\
no se divertían con tu ruina, como los olímpicos.
\end{versopropio}

Odiseo no solo está desengañado de los dioses. Con el tiempo, se ha vuelto ateo y su falta de fe le ha revelado también cuál es la verdadera utilidad de los olímpicos.

\begin{versopropio}
Hace ya mucho tiempo\\
que dejaste de creer en esa gente\\
supuestamente divina, que bebía néctar y daba fiestas\\
en lo alto del monte Olimpo.
\stanzabreak
Lo más probable es que ninguno de tus compañeros,\\
los pocos viejos supervivientes aún en pie,\\
creyera en ellos. Pero todos fingís\\
una devoción apropiada para mantener el statu quo.
\stanzabreak
Y es que, si los dioses no son reales,\\
¿por qué habría la gente de tomarse en serio a los reyes?\\
¿Por qué sufrir a tiranos como vosotros,\\
si están tan vacíos de sentido como las deidades?
\stanzabreak
Así que seguís con las quemas de muslos,\\
y las libaciones, y las plegarias, diligentemente conducidas\\
por esa otra especie de ratas, los sacerdotes,\\
siempre listos para comerse tu comida y beberse tu vino.
\stanzabreak
A cambio, por supuesto, de profecías\\
que resultan coincidir con tus propios planes,\\
y señales que confirman tus antojos,\\
y una bendición general de tus caprichos.
\stanzabreak
Así que, es verdad, los dioses no existen,\\
ni son necesarios, después de todo:\\
los hombres pueden encargarse de abusar,\\
y esclavizar, y matarse unos a otros.
\end{versopropio}

¿Quién, entre las lectoras de cierta edad, entre los lectores que ya peinan canas, no recuerda esos reencuentros de la promoción del instituto, o de la facultad, esas fiestas bienintencionadas para celebrar los veinticinco, o los treinta años del nosotros, los de entonces? ¿Y quién no ha caído en la cuenta de que nosotros, los de entonces, ya no somos los mismos? ¿Quién no se ha asustado con las arrugas de la más guapa de la clase, la barriga cervecera del galán de antaño, los ojos cansados y miopes del otrora poeta, la voz áspera del que cantaba como Orfeo? ¿Quién no ha reconocido su propia decadencia en la de aquellos amigos de otros tiempos?

\begin{versopropio}
Los que todavía se reúnen cada año\\
fingen ser hermanos de sangre,\\
y es cierto, hay calor entre vosotros,\\
algo parecido a la amistad.
\stanzabreak
Ah, pero con qué avidez\\
calibra cada invitado el estado del otro,\\
mide la piel arruinada,\\
el pelo menguante, la espalda rota.
\stanzabreak
Con qué precisión evaluáis\\
los ojos apagados, la mano temblorosa,\\
los párpados caídos, el temblor en la voz,\\
todas las señales que gritan vejez.
\stanzabreak
Este ha engordado demasiado, este está demasiado flaco,\\
Idomeneo se pasa el día durmiendo, Filoctetes siempre está borracho,\\
Néstor apenas puede hablar, la mayoría\\
está en peor forma que tú.
\stanzabreak
Y te preguntas si el único sentido de esas reuniones\\
es reinventar el pasado,\\
espantar la soledad que avanza,\\
y contar los cuerpos supervivientes.
\end{versopropio}

Finalmente, Odiseo reflexiona sobre su amigo Menelao:

\begin{versopropio}
Menelao se conserva bastante bien,\\
puede argumentarse que él y tú\\
sois los menos dañados por el tiempo,\\
sobre todo porque ambos seguís lúcidos.
\stanzabreak
Y él es un verdadero amigo,\\
hay algo de verdad entre las muchas mentiras,\\
fuisteis hermanos de armas,\\
y más de una vez os salvasteis la vida el uno al otro.
\stanzabreak
Así que el vínculo es real,\\
tan real, de hecho, que solo Telémaco\\
es más querido para tu corazón,\\
sin contar, por supuesto, a Circe.
\stanzabreak
Pero tú no le mencionarás la hechicera\\
a tu amigo, como él no mencionará a Helena,\\
los dos habéis aprendido\\
a respetar el corazón roto del otro.
\stanzabreak
No es que seáis adolescentes enamorados,\\
las penas de los jóvenes se arreglan tan fácilmente\\
como sus cuerpos fuertes. Pero el amor perdido de un viejo\\
es el último dolor que llevará a la tumba.
\end{versopropio}

En el siguiente fragmento, Helena ya no está presente en la corte del viejo Menelao:

\begin{versopropio}
Y así, Menelao no habla de Helena,\\
hace ya tantos años que se fue,\\
de hecho, desapareció poco después\\
de la visita de Telémaco a Esparta.
\stanzabreak
Tu amigo presumió tanto ante tu hijo,\\
se aseguró de que viera a la legendaria Helena en todo su esplendor,\\
para que se lo contara a su padre Odiseo,\\
por si acaso había sobrevivido.
\stanzabreak
Y en efecto, Telémaco dijo\\
que nunca había visto a nadie más hermoso,\\
ni más distante. Mirando su fría majestad,\\
se convenció de su naturaleza divina.
\stanzabreak
Así que cuando Helena se desvaneció, el relato oficial\\
habló de una santa ascensión al Elíseo,\\
no se había fugado otra vez, no,\\
sino que los dioses la habían reclamado.
\end{versopropio}

Los versos siguen la narrativa de la \emph{Odisea} y a la vez la subvierten. Siguen así:

\begin{versopropio}
Nadie cuestionó la historia,\\
como nadie cuestiona tu aventura con el cíclope,\\
los cuentos se inventaron precisamente para esto,\\
para encontrar sentido donde solo hay ruido.
\stanzabreak
Ruido y furia, como las hordas que condujiste,\\
rabia, y odio y lujuria, bombeando en vuestros corazones,\\
bronce, y sangre y fuego,\\
para ser disfrazados de batallas épicas.
\stanzabreak
Así que fuisteis todos a Troya,\\
a vengar el rapto de la esposa de Menelao,\\
y luchasteis diez años, y quemasteis la ciudad,\\
y devolvisteis el trofeo al palacio.
\stanzabreak
Pero la segunda vez que ella desapareció,\\
no había ciudad que saquear, ni unidad entre los argivos,\\
y demasiados viejos al mando,\\
así que invocar a los dioses resultó más apropiado.
\end{versopropio}

Odiseo reflexiona sobre el primer rescate de Helena, más motivado quizás por la codicia del botín que por la necesidad de recuperar a la adúltera, y muestra la futilidad de toda la empresa. La primera desaparición da la excusa para una guerra que, seguramente, habría ocurrido de todas maneras; cuando llega la segunda, no pasa nada.

\begin{versopropio}
Lo asombroso\\
es que Menelao la amaba.\\
La amó mucho más después de Troya\\
de lo que la había amado nunca antes.
\stanzabreak
Después de la boda,\\
Helena era solo un trofeo, y que el príncipe,\\
tan joven y furioso, se llevara a la cama a otras mujeres\\
era solo buen deporte, legal, para un hombre.
\stanzabreak
El supuesto secuestro lo enfureció,\\
quizás porque sabía, como todo el mundo,\\
que nadie había obligado a Helena a saltar\\
a la nave de Paris, ni a la cama de Paris.
\stanzabreak
Cuentan los relatos que en Troya\\
Helena desnudó el pecho ante su esposo,\\
lo invitó a golpear, y Menelao\\
dejó caer la espada y lloró.
\stanzabreak
Y por una vez estás convencido\\
de que esta historia sí es cierta:\\
Menelao había necesitado un rechazo y una guerra\\
para entender que estaba enamorado.
\end{versopropio}

El amor a Helena llega tarde y mal y es por ello mucho más trágico. La guerra deja de ser una excusa y se convierte en una justificación y más tarde en un absurdo.

\begin{versopropio}
Pero entonces cometió el error clásico:\\
asumió que su amor recién descubierto,\\
su deseo ferviente, su admiración por ella,\\
le ganarían el corazón.
\stanzabreak
Se equivocaba. Helena, por supuesto,\\
se comportaría como la noble dama que era,\\
y se ocuparía de la casa y del lecho,\\
guardándose sus cartas.
\stanzabreak
Pero fue su esposo quien mató a Deífobo,\\
y tú la viste llorar con el alma rota sobre su cuerpo,\\
qué ironía, pensaste. Ni el valiente Menelao,\\
ni el cobarde Paris, sino aquel otro,
\stanzabreak
príncipe discreto y valeroso,\\
casi tan bravo y honesto como Héctor.\\
A él, y solo a él,\\
le dio Helena su corazón.
\end{versopropio}

Helena, como cualquier personaje literario, admite infinitas lecturas. La mía aquí le hace elegir no al hermoso ---y cobarde--- Paris, ni al aguerrido ---e inútilmente enamorado--- Menelao, sino al más humilde de los héroes, el callado, gentil Deífobo. Toda historia de amor que se precie es una tragedia y la de Helena y Menelao no es una excepción. Mientras el príncipe aqueo bebe los vientos por su esposa, ella trama, fría y pacientemente, su venganza:

\begin{versopropio}
Y por él, y solo por él,\\
tramó su venganza.\\
No eligió el cuchillo, como su hermana,\\
ni el veneno que tan bien conocía.
\stanzabreak
No: ella quería herir más hondo,\\
así que trabajó duro hasta que Menelao\\
solo podía respirar el aire que ella había respirado,\\
y entonces lo apuñaló por la espalda.
\stanzabreak
No hizo falta bronce, ni siquiera las nuevas dagas de hierro.\\
Ella, simplemente, se fue;\\
dicen algunos que un troyano llamado Eneas\\
organizó la expedición de rescate.
\end{versopropio}

Así que en la vuelta final de tuerca, Helena huye con Eneas, rompiendo, esta vez sí, el corazón de Menelao:

\begin{versopropio}
Menelao no habla de ella,\\
como tú no hablas de Circe,\\
cada año, los dos os reunís\\
en Esparta y bebéis vino espeso,
\stanzabreak
y os emborracháis, y lloráis por el amigo perdido,\\
y por el amor perdido\\
que nunca merecisteis.
\end{versopropio}

A eso me refiero cuando hablo de reinventar un personaje como Helena. Nadie pide fidelidad a Homero, pero sí estar a la altura de sus personajes.
